\scanpage{41}

Démonstration. On \deleted{\unsure{considère|conditions}} sait que $f_*(U_{ij}/T)$ existe \unsure{aussi}, il s'envoie par imm. ouverte dans \deleted{les} $f_*(U_i)$ \uline{et} $f_*(U_j)$. Les conditions de cohérence sont \unsure{manifestemt} vérifiées, d'où un préschéma $\underline H$ sur $S$, \deleted{\unsure{on a une}} réunion d'ouverts $H_i\simeq f_*(U_i/T)$.

On déduit par \unsure{recollmt|recollement} un $T$-hom $H\times_S\deleted{\unsure{S}}T\to X$, \inserted{\deleted{\unsure{induisant|induit|universels}} $H_i\times T\to U_i$.} Je dis que ce dernier définit $T$ comme un \deleted{\unsure{Hom}} $f_*(X/T)$.

% slanted
\margin{On vérifie sans difficulté que $H_i$ est le plus gr\textsuperscript{d} ss-ouv de $H$ \deleted{\unsure{contenu}} dont l'image inverse dans $H\times_S T$ est dans l'image réciproque de $U_i$}

Soit $S'\to\deleted{\unsure{H|S}}$ \deleted{\unsure{un schéma, prouvons}} \inserted{sur $S$, il faut \unsure{savoir|prouver}} que $\mathrm{Hom}_S(S',H)\xrightarrow{\sim}\mathrm{Hom}_T(S'\times_S T,X)$.

On peut supposer $S=S'$.

\deleted{La flèche est \uline{injective}, car si on a \deleted{\unsure{une section}}\inserted{deux} $\deleted{\unsure{S\to T|S'\to S}}$, \deleted{\unsure{la s :}} sections \deleted{\unsure{$s',s''$}} $u,v:S\to H$ définissant le $=$ \deleted{application}\inserted{morphisme} $T\to X$, les images inverses $T_i$ par ce dernier \deleted{\unsure{de $U_i$|ouverts}} des $U_i$ \uline{recouvrent} $T$,}

\diagram{\begin{tikzcd}[ampersand replacement=\&,column sep=small,row sep=small]
\deleted{S}\arrow[r,dash]\arrow[d,dash] \& \deleted{H}\arrow[d,dash]\\
\deleted{\unsure{S\times T|S'}}\arrow[r] \& \deleted{\unsure{S|T}}
\end{tikzcd}}

\diagram{\begin{tikzcd}[ampersand replacement=\&,row sep=small,column sep=small]
\& X\arrow[d]\supset U_i \&\\
S \& T\arrow[l] \&\\
H\arrow[u] \& \& H\times_S T=\bigcup H_i\times_S T\arrow[ll]\arrow[ul]\arrow[uul,bend right=18]
\end{tikzcd}}

i.e. il faut prouver $\mathrm{Hom}_S(S,H)\simeq\mathrm{Hom}_T(T,X)$\unsure{.:|.|,} Soit $\varphi$ une section de $X/T$,

% archive sheet number: 40
